
\section{Solution Methodology for LRSP}
\label{SolnMethod}

** different espprc: one is an espprc algorithm, 1 label for 
each facility. Instead of solving pricing problems for each facility, solve 
1 pricing!

The set partitioning formulation AUG-SPP-LRS contains one variable for each possible
pairing. For practical size problems, the number of pairing will be very
large, and, as a result, it is unlikely that we can solve instances that 
explicitly include all feasible pairings. Instead, we have developed a branch-and-price
algorithm, in which we use column generation to solve the linear programming 
relaxations at the nodes of the tree. 

In sections \ref{ColGenAlg}-\ref{Solve-pricing}, we describe the details of the column generation
algorithm, including the solution methods for the pricing problem. In section \ref{BPAlg}, 
we describe the details of the branch-and-price algorithm, including the branching 
rules used. In section \ref{BPAlg-Opt}, we discuss the conditions under which the
nodes of the branch-and-price tree can be fathomed.   

\subsection {Column Generation}
\label{ColGenAlg}
Column generation algorithms have been applied to a wide variety of large scale 
optimization problems, such as vehicle routing (\citet{Agarwal}, \citet{Desrochers2}), 
airline crew scheduling (\citet{Anbil}, \citet{Vance}), location and routing (\citet{Berger05}), 
cutting stock (\citet{Vance2}), lot sizing and machine scheduling (\citet{Vanderbeck}). 
\citet{Lubbecke} list other problems to which column generation has been applied.

Column generation algorithms are applied to models that contain too many columns
to be handled explicitly.
By solving a sequence of linear programs and using dual information, the algorithm 
explicitly considers all of the columns but generates only those that may improve 
the objective function. We use column generation (CG) to solve appropriate relaxations
of the AUG-SPP-LRS formulation at all nodes of the branch-and-price tree. In this section,
we describe details of the column generation algorithm for the root node of the tree.

To begin a CG procedure, we need an initial set of columns that contains a feasible 
solution. We call the problem that contains only a subset of the columns of the original 
problem the {\bf restricted master problem (RMP)}. From the solution of the RMP, we use 
the associated dual variable information to identify new columns to add to the RMP. We 
solve an optimization problem called the {\bf pricing problem} to identify new columns.

A {\bf column generation algorithm} for solving the LP relaxation of AUG-SPP-LRS 
includes the following four steps: 

\begin{enumerate}
\item[0.] Generate an initial set of columns that contains a feasile solution.
\item[1.] Solve the RMP and obtain dual variable information.
\item[2.] Using the dual variables, form a pricing problem that seeks to identify 
new pairings (columns) and their associated reduced costs. 
\item[3.] Solve the pricing problem. If the solution specifies a column with negative
reduced cost, add the new column to the RMP and go to step 1. Otherwise, the 
solution found in step 1 is an optimal solution for the LP relaxation of AUG-SPP-LRS.
\end {enumerate}

The challenge in applying this algorithm to solve the LP relaxation of AUG-SPP-LRS 
is in steps two and three. We need to define a pricing problem that identifies
feasible pairings with negative reduced cost and construct an algorithm for
solving the problem efficiently.
  
\subsection {Formulating the Pricing Problem}
\label{form-pricing-prob}
The objective of the pricing problem is to generate columns with 
negative reduced cost to add to the RMP. We know that, in an
LP that is a minimization problem, the reduced costs of the variables
are zero or positive in an optimal solution. Therefore, if the pricing problem does not
identify any columns with negative reduced cost, then the current  
solution of the RMP is optimal. For the LRS formulation, columns represent feasible 
pairings.

To formulate the pricing problem, we first must associate dual variables  
with the constraints. Let $\pi_{i}$ be the dual variable associated with the
partitioning constraint (\ref{spp_con1}) for customer i, $\mu_{j}$ be 
the dual variable associated with the capacity constraint (\ref{spp_con2}) 
for facility j and $\sigma_{ji}$ be the dual variable for the linking 
constraint (\ref{relation}) for facility j and customer i. Then, for the 
variable $Z_{jp}$ for $j \in M$ and $p \in P_{j}$, the reduced cost $\hat{C}_{jp}$
can be written as:   
\begin{eqnarray}
\label{red-cost}
\hat{C}_{jp}=C_{jp}-\sum_{i\in N}a_{ipj}\cdot \pi_{i}+\sum_{i \in N}a_{ipj} \cdot d_{i} \cdot \mu_{j}+\sum_{i \in N}a_{ipj} \cdot \sigma_{ji}
\end{eqnarray} 

Note that the reduced cost of a pairing (\ref{red-cost}) does not include a dual 
variable associated with the constraint (\ref{lb-open-fac}), since the coefficients of 
the $Z_{jp}$ variables in constraint (\ref{lb-open-fac}) are zero.

The objective of the pricing problem is to identify a column (pairing) 
that has the minimum $\hat{C}_{jp}$ value. Remember that a pairing is a set 
of vehicle routes that satisfy the following constraints:

\begin{description}
\item [C1.] Total demand of each route must not exceed the vehicle capacity.
\item [C2.] Total travel time of the routes in the pairing  must not exceed the 
time limit.
\item [C3.] All of the routes in the pairing must start and end at the same facility.
\item [C4.] Each customer node can be visited at most once.
\end{description}

Since every route in a pairing must start and end at the same facility, 
each facility has its own set of feasible pairings, which allows us to decompose
the pricing problem into $|M|$ independent pricing problems. Therefore, we solve one 
pricing problem for each facility to search for columns with negative reduced costs. 

Recall that in equation (\ref{red-cost}), $C_{jp}$ is the cost of pairing $p$ associated with 
facility $j$. The cost $C_{jp}$ is the sum of the fixed cost of the vehicle and the variable 
costs of the routes that are included in the pairing. The variable cost of a route is equal to 
the sum of the costs of arcs in the route. If we can assign the components of expression  
$-\sum_{i\in N}a_{ipj} \cdot \pi_{i}+\sum_{i \in N}a_{ipj}\cdot d_{i}\cdot \mu_{j}+\sum_{i \in N}a_{ipj} \cdot \sigma_{ji}$ to arcs, then the reduced cost of a pairing, $\hat{C}_{jp}$, can be expressed as the 
sum of the costs of the arcs included in the pairing plus the fixed cost of the vehicle. 
Although each dual variable is associated with a customer node $i$ and/or a facility node $j$ 
instead of an arc, later  we will show how to transfer these node related dual variables 
to the arcs. In this way, the reduced cost of a pairing can be computed as the sum of the 
reduced cost of the arcs in the pairing.     

\subsubsection{Modified ESPPRC}
The pricing problem, which seeks to find a pairing of minimum reduced cost $\hat{C}_{jp}$ 
subject to constraints (C1), (C2), (C3) and (C4) can be defined as an elementary shortest 
path problem with resource constraints (ESPPRC) on a modified network. In an {\bf ESPPRC}, 
the objective is to find a shortest path between the source and the sink such that each node 
is visited at most once and the resource constraints are not violated. To properly represent
our pricing problem as an ESPPRC, we need to make one adjustment; while in an ordinary 
elementary shortest path problem, each node can be visited at most once, in our case, 
because we are looking for a set of routes, we allow the sink node to be visited more than 
once. In our modified ESPPRC, the sink node represents the end of one route and maybe the 
beginning of another. In an ordinary ESPPRC, a path from the source to the sink corresponds 
to one route, however, a solution of our modified ESPPRC may correspond to multiple routes. 
Another reason for multiple visits to the sink is to reset the remaining capacity of the vehicle 
before the new route. In a pairing, the vehicle starts a route with a full truck 
load and delivers the required demand to all of the customers on the route. The vehicle returns 
to the depot after the last customer in the route and resets its capacity to full truck load for 
the next route. In the pricing problem, the truck capacity is one of the resources. Whenever the 
path visits the sink in the modified network, the vehicle capacity resource consumption is updated and 
set to full truck load. 

\subsubsection{Constructing Network for Modified ESPPRC}
%The definition of a pairing allows us to define the pricing problem as a 
%network problem in which we determine a minimum cost path from a 
%source to a sink. 
To define the pricing problem for facility j, we construct a network with $|N|+2$ nodes, 
$|N|$ nodes for the customers, one for facility $j$ as a source node and a copy of facility 
$j$ as a sink node. Each customer node has a demand equal to its demand in the original problem. 
Let $A^j$ be the arc set of the network for facility j. Then $A^j$ consists of four sets of arcs: 
\begin{eqnarray*}
A^{j}_{s,c} & = &\mbox{set of arcs from the source to customer nodes} \\ 
A^{j}_{c,si} & = &\mbox{set of arcs from customer nodes to the sink} \\
A^{j}_{si,c} & = &\mbox{set of arcs from the sink to customer nodes} \\
A^{j}_{c,c}  & = &\mbox{set of arcs between pairs of customer nodes}
\end{eqnarray*}

Note that there are no arcs from customer nodes back to the source but there are arcs from the sink
to customer nodes. The arcs from the sink to customer nodes represent the beginning of a 
new route whereas the arcs from the customer nodes to the sink represent the end of a route.
Figure \ref{fig:exp_network} shows an example network and a path on the network. 
The solution of the  pricing problem is a path from the source to the sink on the constructed 
network. If the solution visits the sink  more than once, then the solution corresponds to 
a pairing that includes more than one route. 

\begin{figure}[h]
  \centering
  \includegraphics[scale=0.8]{./figs/NetworkexpG.eps}
  \caption{Example of a network and a path from source to sink (dashed line)} \label{fig:exp_network}
\end{figure}


%A visit to the sink resets the remaining capacity of the vehicle to a full truckload. 

%In a pairing, the vehicle starts a route with a full truck 
%load and delivers the required demand to all of the customers on the route. The vehicle returns 
%to the depot after the last customer in the route and resets its capacity to full truck load for 
%the next route. In the pricing problem, the truck capacity is one of the resources. Whenever the 
%path visits the sink in the modified network, the vehicle capacity resource consumption is updated and 
%set to full truck load. 
       
%$A^{j}_{si,c}$, arcs from sink to customers and 
%$A^{j}_{c,c}$, arcs between any two customers. There are no arcs from customer 
%nodes to the source. 

To find a pairing of minimum reduced cost using the constructed network, we need to find an 
appropriate way to incorporate the dual variables into the costs of the arcs in  $A^j$. 
Remember that $t_{ik}$ is the travel time from location $i$ to location $k$ and $d_i$ is the 
demand of customer $i$ and let $c_{ik}$ be the operating cost of traveling from location 
$i$ to location $k$ in the LRS problem. To define the length and the cost of arcs on the 
constructed network, let $t^{'}_{ik}$ be the travel time on arc $(i,k)$ and $c^{'}_{ik}$ 
be the reduced cost of traveling arc $(i,k)$. Then, for the network associated with 
facility j:  

%Since in the pricing problem, we aim to minimize the reduced cost of the variable 
%$Z_{jp}$ (\ref{red-cost}), if we want to solve the pricing problem as a network problem

\begin {itemize}
\item {\em Case 1:} If $(i,k)$ is an arc from source or sink to a customer node, 

  e.g., $(i,k) \in (A^{j}_{s,c} \cup A^{j}_{si,c})$, then 
  $t^{'}_{ik}=t_{jk}$, $c^{'}_{ik}=c_{jk}-\pi_{k}+d_k \cdot \mu_j+\sigma_{jk}$ 
\item {\em Case 2:} If $(i,k)$ is an arc between two customer nodes, 

  e.g., $(i,k) \in A^{j}_{c,c}$, 
  then   $t^{'}_{ik}=t_{ik}$, $c^{'}_{ik}=c_{ik}-\pi_{k}+d_k \cdot \mu_j+\sigma_{jk}$ 
\item {\em Case 3:} If $(i,k)$ is an arc from a customer node to sink, 

  e.g., $(i,k) \in A^{j}_{c,si}$, then $t^{'}_{ik}=t_{ij}$, $c^{'}_{ik}=c_{ij}$
\end{itemize} 

We transfer the dual variables $\pi_i$ and $\sigma_{ji}$ associated with customer 
$i$ to the incoming arcs of customer node $i$. As in the reduced cost calculation
in (\ref{red-cost}), the dual variable for facility $j$, $\mu_j$, is multiplied with
demand value $d_i$ for customer $i$ that is in the pairing. Therefore, the cost 
$d_i \cdot \mu_j$ is also carried on to the incoming arcs of customer node $i$. Since facility
$j$ does not have demand, the dual variable $\mu_j$ does not affect the arcs 
into the sink node.  

To illustrate why a path through the modified network is equivalent to a pairing, we consider
the following example. Let $N=\{1,2,3,4,5,6,7,8,9,10\}$ be the set of demand nodes.  Then the 
set of all nodes in the network is $I=\{source,1,2,3,4,5,6,7,8,9,10,sink\}$. Consider the path
(source-2-4-sink-5-6-7-sink-9-10-sink) in the modified network. This path in the modified network
is equivalent to a pairing with three routes in our pricing problem. The component routes are 
($facility_j$-2-4-$facility_j$), ($facility_j$-5-6-7-$facility_j$) and ($facility_j$-9-10-$facility_j$). 
The pairing is feasible if each route's total demand does not exceed the vehicle capacity and if the 
sum of the travel times of the paths does not exceed the time limit. From the arc costs  specified 
in Cases 1 through 3, the reduced cost of the path in the modified network is:
\begin{eqnarray*}
\mbox{Total cost of the path} &=& VF+(c^{'}_{j,2}+c^{'}_{2,4}+c^{'}_{4,j}+c^{'}_{j,5}+c^{'}_{5,6}+c^{'}_{6,7}+ c^{'}_{7,j}+c^{'}_{j,9}+c^{'}_{9,10}+c^{'}_{10,j})\\
                              &= &VF+(c_{j,2}-\pi_2+d_2\mu_j+\sigma_{j,2})+(c_{2,4}-\pi_4+d_4\mu_j+\sigma_{j,4}) \\
                              &+ &(c_{4,j})+(c_{j,5}-\pi_5+d_5\mu_j+\sigma_{j,5})+(c_{5,6}-\pi_6+d_6\mu_j+\sigma_{j,6})\\
                              &+ &(c_{6,7}-\pi_7+d_7\mu_j+\sigma_{j,7})+(c_{7,j})+(c_{j,9}-\pi_9+d_9\mu_j+\sigma_{j,9})\\
                              &+ &(c_{9,10}-\pi_{10}+d_{10}\mu_j+\sigma_{j,10})+(c_{10,j})
\end{eqnarray*} 
Grouping terms differently and noting that the column vector representing this pairing is 
$a_{jp}=[0 1 0 1 1 1 1 0 1 1]$ and the cost of a pairing is the sum of the arc costs plus the vehicle 
fixed cost, we obtain the following expression, which is just the reduced cost of $Z_{jp}$ as calculated
via equation (\ref{red-cost}).
\begin{eqnarray*}
\hat{C}_{jp} &=& VF + (c_{j,2}+c_{2,4}+c_{4,j}+c_{j,5}+c_{5,6}+c_{6,7}+ c_{7,j}+c_{j,9}+c_{9,10}+c_{10,j})\\
             &-& (\pi_2+\pi_4+\pi_5+\pi_6+\pi_7+\pi_9+\pi_{10})+\mu_j \cdot (d_2+d_4+d_5+d_6+d_7+d_9+d_{10})\\
             &+& (\sigma_{j,2}+\sigma_{j,4}+\sigma_{j,5}+\sigma_{j,6}+\sigma_{j,7}+\sigma_{j,9}+\sigma_{j,10})\\
             &=& C_{jp}-\sum_{i\in N}a_{ipj} \cdot \pi_{i}+\sum_{i \in N}a_{ipj} \cdot d_{i}.\mu_{j}+\sum_{i \in N}a_{ipj} \cdot \sigma_{ji}
\end{eqnarray*} 
Therefore, the reduced cost of a pairing (sequence of paths) in the original network is equal to the cost
of a path in the modified network.


\subsection{Solving the Pricing Problem}
\label{Solve-pricing}
In this section, we describe approaches for solving two versions of the pricing problem, a relaxation 
and the original. There are both positive and negative aspects associated with each version. 

{\bf Solving the Relaxation of ESPPRC}

One approach to find negative reduced cost pairings is to solve a relaxation of the pricing problem in which we 
allow customer nodes to be visited more than once. To incorporate this change into the AUG-SPP-LRS formulation, we change the 
definition of the parameter $a_{ijp}$ to represent the number of visits to customer node $i$ in the pairing $p$ 
associated with the facility $j$. The advantage of the relaxed version is that it can be solved using a dynamic 
programming algorithm. The complexity of the algorithm 
is pseudopolynomial depending on the number of resource constraints (\cite{Beasley}). In our 
case, we have two resources: the vehicle capacity ($VCap$) and the time limit ($TL$). Therefore, the algorithm
has $O(|TL| |VCap| |N|^2)$ complexity. A disadvantage of this approach is that, because one  set of constraints 
is relaxed, the LP relaxation bound may be weakened. That is,  when the CG algorithm
terminates, the LP relaxation solution may contain walk columns, in which case the objective function is only 
a lower bound on the optimal LP relaxation bound. In section \ref{comp-espprc-dp}, we demonstrate empirically
how the bound of this relaxation compares to the optimal LP relaxation bound. 

To solve the relaxed version, we have developed a standard forward dynamic programming algorithm, which finds 
minimum cost walks from the source 
to the sink considering the vehicle capacity and the time limit constraints. \citet{Beasley} give a formal 
description of such a dynamic programming algorithm.  
We have three state variables: the current node, the accumulated time and the current vehicle load.
The state variable representing the accumulated time increases  whenever we move
from one node to another. However, the current vehicle load may increase or decrease
depending on the node to which we arrive. Since we assume that the problem is a delivery problem,  
the vehicle load decreases when we arrive to a customer node.  The vehicle load increases and
is set to the full truckload when we arrive to the sink node. To reset the vehicle load to the full truck 
load, we set the demand of the sink to $-VCap$. If we let $(i, T^t, T^v)$ be the state variables where 
$i$ represent the current node, $T^t$ represents
the accumulated time from the source to the node $i$ and  $T^v$ represents the current vehicle load, then
$c(i, T^t, T^v)$ is defined as the state cost. The forward recursive function is:
\begin{eqnarray*}
&c (source,0,0)    =   0 \hspace{3.02cm} & \\
&c (j, T^t, T^v)   =   \min_{i} \{ c'_{ij}+c(i,T_i^t,T_i^v)| &T_i^t+t'_{ij}= T^t \land T^t+t'_{j,sink}\leq TL \land \\
                        & &\min(VCap,T_i^v-d_j)=T^v \land T^v \geq 0\}
\end{eqnarray*}
The first and the third constraints in the definition of the recursive function update the state variables. 
The third constraint guarantees that the vehicle load is reduced when $j$ is a customer node, and
the vehicle load is set to vehicle capacity if $j$ is the sink. The second and 
the last constraints are used to check the feasibility of the state. 


{\bf Solving the Exact ESPPRC }

Solving the original pricing problem in which a customer may be visited at most once is more difficult
than solving the relaxation because we must keep track of whether a node has been visited
or not. \citet{Beasley} describe how to 
find elementary paths in a graph by adding an extra binary resource for each node showing 
whether the node is visited or not.
For  shortest path problems with resource constraints (SPPRC), \citet{Desrochers}  
developed a label correcting algorithm that is based on dynamic programming and that has  
pseudopolynomial complexity. The algorithm has been shown to be  successful when there 
are tight resource constraints, but it becomes more time consuming as the number of resources increases. 
As the algorithm builds partial paths, it assigns a label to each 
indicating the resource consumption. To eliminate the problem of an exponentially 
increasing number of labels, the algorithm applies dominance rules to delete dominated labels. 
\citet{Feillet} adapt \citeauthor{Desrochers}'s label correcting 
algorithm to find elementary paths and strengthen the algorithm by developing additional domination 
rules for the labels. 

We have adapted \citet{Feillet}'s label correcting algorithm to our pricing problem to solve the  
ESPPRC described in section \ref{form-pricing-prob}. In the algorithm, each path from the source to any node 
in the network is assigned a 
label. Let $X_{so,i}$ be a path from the source node to node $i$. The label $L^{X_{so,i}}$ 
that corresponds to path $X_{so,i}$, is associated with a state $(T_i^{v},T_i^{t},s_i,V_i^1,..,V_i^N)$  
that keeps the information about the path. In the state,  $T_i^{v}$ is the current vehicle load,
$T_i^{t}$ is the time consumption of the path and the vector $(V_i^1,..,V_i^N)$ gives
the sequence of nodes in the path as well as the unreachable nodes for the path. 
The elements of the vector, $(V_i^1,..,V_i^N)$ are defined as:
\begin{eqnarray*}
V_i^k & = & \left\{ \begin{array}{llll}
                 -1 & \mbox{if node $k$ is {\bf unreachable} from node $i$ for the path because of the}\\
                    & \mbox{ resource limits, }\\
                  0 & \mbox{if node $k$ is not in the path and it is {\bf reachable} from node $i$ with the }\\
                    & \mbox{current resources,}\\
                 n>0 & \mbox{if node $k$ is {\bf unreachable} from node $i$ since it is in the path and it is} \\
		     & \mbox{the $n^{th}$ node.} 
                \end{array} 
              \right. 
\end{eqnarray*}
The state variable $s_i$ keeps the number of unreachable nodes. In the label $L^{X_{so,i}}$,  node $k \in N$ is 
{\em unreachable} from node $i$ if it is already in the path or if there is not 
enough of a resource to extend the path (label) from node $i$ to node $k$. For example, $k$ is unreachable
if the vehicle load is not enough to serve the demand of customer $k$ ($T_i^v-d_k<0$) and it is not possible to 
travel to node $k$ and then to the sink (to finish the path) within the time limit ($T_i^{t}+t'_{ik}+t'_{k,sink}>TL$).   

Note that the dimension of the vector $(V_i^1,..,V_i^N)$ is $|N|$, which means that we have an element for each
customer node in the network. We do not need an element for the source, because every path starts from the source 
and can never return to it. Similarly, we do not need an element for the sink, because the sink is always reachable 
for a path by construction. Extending a path to the sink is always resource feasible: (1) at the sink, the vehicle 
load is reset to the capacity of the vehicle and (2) when constructing a path, the time constraint  
is checked to ensure that the path is extendable to the sink; otherwise it is not created.    
The label $L^{X_{so,i}}$ for the path $X_{so,i}$ also has an associated cost parameter $C_{X_{so,i}}$ that keeps
the reduced cost of the path.   

The label correcting algorithm for solving the ESPPRC includes three main parts: 
\begin{enumerate}
\item {\bf Initialization step} To begin the algorithm, a list of labels is created 
  for each node to keep the created labels for the node. The list of  labels for all of the 
  nodes except the source node is initially empty. The source has one label with 
  all of the state variables and the reduced cost set to zero.
\item {\bf Extension step} For each node $i$, the algorithm examines its list of labels
  and tries to extend them to all of the successors of node $i$. A label for node $i$ can be
  extended to node $k$ if one of the following is true: (i) the node $k$ is the sink or (ii) 
  the node $k$ is reachable from node $i$ ($V_i^k=0$) and the extended label is resource 
  feasible ($T_i^v-d_k\geq 0$ and $T_i^{t}+t'_{ik}+t'_{k,sink}\leq TL$). A new label is
  created for node $k$ by extending the label of node $i$, that is, by updating the resource 
  consumptions, the vector $(V_i^1,..,V_i^N)$ and the reduced cost. 
\item {\bf Elimination step } After extending all possible labels of all nodes to 
  construct new labels for their successors, we combine these new labels  
  with their previously created label lists. Let node $k$ be a node with new labels. 
  %After constructing new labels for node $ k$ in the network 
  %from the extendable labels of all nodes for which the node $k$ is a successor, 
  %we combine these new labels of node $k$ with the label list of $k$ that includes previously 
  %created labels. 
  We combine the label list of $k$ with the list of new labels created for $k$. By comparing each 
  label with the others in the list, we determine the {\em non-dominated labels}. 
  Let $L_k^1$ and $L_k^{2}$ be two labels in the list.  Let 
  $(T_k^{1v},T_k^{1t},s_k^1,V_k^{11},..,V_k^{1N})$ and 
  $(T_k^{2v},T_k^{2t},s_k^2,V_k^{21},..,V_k^{2N})$  be the corresponding state variables 
  and $C^1$ and $C^{2}$ be the reduced costs of the labels. Label $L_k^1$ {\em  dominates} label
  $L_k^{2}$, if all of the following conditions are satisfied:
  \begin{itemize}
  \item The current resource consumptions of $L_k^1$ are less than or equal to the resource consumptions
    of $L_k^{2}$, i.e., $T_k^{1v}\leq T_k^{2v}$ and $T_k^{1t}\leq T_k^{2t}$.
  \item The number of unreachable nodes in $L_k^1$ is less than or equal to the number in $L_k^{2}$, i.e.,
    $s_k^1\leq s_k^2$.
  \item $V_k^{1p}\preccurlyeq V_k^{2p} \ \ \forall p=1..N$, where the relation
    $V_k^{1p}\preccurlyeq V_k^{2p}$ is true if $V_k^{1p}=0$ (p is reachable from k) in $L_k^{1}$ 
    or $p$ is unreachable from $k$ in both labels.
  \item The reduced cost of $L_k^1$ is less than or equal to the reduced cost of $L_k^{2}$, i.e.,
    $C^1 \leq C^{2}$.
  \item The labels are not equal, i.e.,

    $(C^1,T_k^{1v},T_k^{1t},s_k^1,V_k^{11},..,V_k^{1N})$ $\neq$ 
    $(C^2,T_k^{2v},T_k^{2t},s_k^2,V_k^{21},..,V_k^{2N})$.
  \end{itemize}
   All dominated labels are deleted from the label lists of the nodes. If any label list is changed
   after the elimination step, i.e., new labels are added, then the algorithm returns to the extension 
   step. When all the labels are extended and no new labels are added to the lists of nodes, the 
   algorithm terminates. 
\end{enumerate}
When the algorithm terminates, the labels in the list for the sink node corresponds to feasible pairings.
The labels with negative reduced cost are candidate pairings to be added to the RMP. If the algorithm 
concludes with no negative reduced cost labels on the sink node, then we can 
conclude that the current RMP is optimal. 

{\bf Additional Domination Rule}

In addition to \citet{Feillet}'s domination rules, we have developed
a domination rule which reduces the symmetry. Since in the pricing problem we are 
looking for a set routes instead of one route  - the algorithm is
not elementary with respect to the sink node - the state space of the problem is
larger than an ordinary ESPPRC. Although the ordering of the routes within a pairing
is arbitrary, a path from source to sink in the constructed network corresponds to 
a pairing including an ordered set of vehicle routes. To eliminate generation of all 
possible ordered sets, we introduce a new domination rule in which we force a specific 
order for routes in the pairing, which decreases the state space of the algorithm 
significantly. The following rule is added to the step 3 of the label correcting 
algorithm.
\begin{itemize}
\item [{\bf Rule}] The index of the first customer node of each route should be less 
than the index of the first customer node in the following route.
\end{itemize}
The modified algorithm only generates paths 
from source to sink in which the first node of each route is less than the first node 
of the following route. 
For example, let (source-1-2-sink-3-9-4-sink-5-7-sink) be a 
feasible path visiting customer nodes in the set \{1,2,3,4,5,7,9\}. We can view this 
pairing also as a combination of the following three routes: facility-1-2-facility, 
facility-3-9-4-facility and facility-5-7-facility. The modified algorithm does not 
generate any of the following paths: (source-1-2-sink-5-7-sink-3-9-4-sink), 
(source-5-7-sink-1-2-sink-3-9-4-sink), (source-5-7-sink-3-9-4-sink-1-2-sink),
 (source-3-9-4-sink-1-2-sink-5-7-sink),  (source-3-9-4-sink-5-7-sink-1-2-sink)
since the first nodes of the routes in these pairings are not in increasing order. 

\zaupdate{ {\bf 1. TWO -STEP ESPPRC ALGORITHM} \\
**************************}


\zaupdate{ {\bf 2. PRICING PROBLEM WITH LABEL LIMIT }\\
*****************************}


\zaupdate{ {\bf 3. BIDIRECTIONAL PRICING PROBLEMS} \\
*********************************}



{\bf Solving the Pricing Problem in Column Generation Algorithm}
 
We use two approaches to find a column with negative reduced cost in our branch-and-price 
algorithm:  an {\em exact approach} and a {\em heuristic approach}. In the 
exact approach, we solve the modified ESPPRC label correcting algorithm and the algorithm 
gives us a column with most negative reduced cost as well as some other columns with negative 
reduced cost. In the heuristic approach, we solve the modified ESPPRC label correcting algorithm
with a label limit  at each node. We set a label limit ($LL$), the maximum number of
labels will be kept at each node during the algorithm. When we update the label list for 
a node, we order the labels in the list in increasing order of reduced cost. We keep at most
the $LL$ number of labels of the ordered label list at each node. If the heuristic ESPPRC algorithm results
columns with negative reduced cost, they are feasible for the original problem and they can
improve the LP relaxation solution, but they are not necessarily the columns with most negative 
reduced cost. 



\subsection{Branch-and-Price Algorithm}
\label{BPAlg}
In this section, we describe our branch-and-price algorithm for the LRS problem.
Branch-and-price algorithms, which combine column generation with branch-and-bound 
algorithms, have been applied to many different problems in recent years, such as 
assignment problems (\citet{Savelsbergh}), airline crew scheduling problems (\citet{Vance}) 
and location routing problems (\citet{Berger05}).

%\citet{Savelsbergh} applies a branch-and-price algorithm to a generalized 
%assignment problem which is a structure that appear in many models, such as vehicle routing, 
%facility location and scheduling. He showed that the algorithm outperforms many of the other 
%optimization algorithms used before. \citet{Vance} formulate an airline crew scheduling 
%problem and describe many details of their implementation of the branch and price algorithm. 
%Although the LP relaxation of their formulation gives a stronger bound than traditional 
%formulations, their LP relaxation is more difficult to solve than the traditional ones. 
%Closely related to our work is the work of \citet{Berger}, in which the authors 
%develop a branch and price algorithm for an uncapacitated location and routing problem.  

%[Although branch-and-price has been applied successfully to some of the 
%large scale optimization problems such as crew scheduling, machine
%scheduling, facility location and routing, since the mid-1990s, there are 
%many difficulties associated with the algorithm and the success of the 
%algorithm highly depends on the structure of the problem. \citet{Barnhart}
%provide an overview of the general methodology and many of the challenges.]

\citet{Barnhart} present a discussion of the general methodology and many of the challenges
of developing a branch-and-price algorithm. The basic idea of a branch-and-price algorithm is to 
solve an appropriate LP at each node of the branch-and-bound tree using column generation. Columns must be 
generated throughout the tree because we might not have all possible columns in the current 
node and branching rules and fixing of variables may change the values of the dual variables. 
Therefore, the reduced costs of the columns may change and there may exist new columns with negative 
reduced costs that are not in the current problem. The main challenge in developing a 
branch-and-price algorithm is developing effective branching rules that do not destroy the 
structure of the pricing problem. Figure \ref{fig:bp_tree}  shows a schematic representation 
of a branch-and-price algorithm.

 
%There are some remarks about branch-and-price algorithm that we have to state.
%In every node of the branch and bound tree, column generation should be used to find the 
%optimal LP solution at the corressponding node. While  pricing new columns in column generation
%algorithm, fixed variables and branching rules must also be considered, which means that the 
%pricing problem must be modified. {\bf Figure 2} shows the schematic representation of the 
%branch-and-price algorithm.

\begin{figure}[h]
  \centering
  \includegraphics[scale=0.5]{figs/figure2.eps}
  %\fbox{ \includegraphics[scale=0.5]{figure2.eps}}
  \caption{Branch-and-Price Tree}\label{fig:bp_tree}
\end{figure}

Our branch-and-price algorithm includes the following steps, which are described in
general terms. %We will provide details of each step later.       

\begin {enumerate}
\item {\bf Solving the Root Node of the Branch-and-Price Tree} 
  \begin {enumerate}
  \item [{\bf Step 1.}] Starting with an initial set of columns, solve the LP relaxation 
    of the AUG-SPP-LRS using a column generation algorithm.
  \item [{\bf Step 2.}] The LP relaxation solution is a lower bound for the problem.  
    If the LP relaxation solution satisfies the integrality constraints, 
    then it is an optimal solution of the original problem and the algorithm stops.  
    Otherwise, go to Step 3. 
  \item [{\bf Step 3.}] Find an integer feasible solution of the RMP with the current set of 
    columns using a branch-and-bound algorithm. Set this solution as the upper bound for the problem. 
    Go to step 4. 
  \end{enumerate}
\item {\bf Branching} 
  \begin{enumerate}
  \item [{\bf Step 4.}] Select one or more fractional variables on which to branch. 
    Create new nodes in the branch-and-price tree. 
  \item [{\bf Step 5.}] Do one of the following: 
    \begin {itemize}
    \item If there is an unprocessed leaf node in the tree, select one to process. Fix the 
      variables  according to the branching 
      rules applied along the path in the tree from the root node to the selected node. Modify the 
      pricing problem to incorporate the branching rules. Go to Step 6.
    \item If all nodes in the tree are marked as processed, then STOP. The current best integer
      feasible solution is an optimal integer solution.
    \end{itemize}
  \item [{\bf Step 6.}] Solve the linear program associated with the current node using 
      a column generation algorithm with the modified pricing problem. 
      \begin {itemize}
      \item If the linear program at the current node is infeasible, then fathom the node,
	mark it as processed and return to Step 5. 
      \item If the LP solution at the current node satisfies the integrality constraints, 
        then it is a feasible solution for the original problem. If the current solution
        is better than the best integer feasible solution, then update the upper bound. 
	Return to Step 5. 
      \item If the LP solution contains fractional variables, go to Step 4. 
    \end{itemize}
  \end {enumerate}
\end{enumerate}

%In the previous sections we have showed that how we formulate our initial 
%pricing problem and how we solve the root node using column generation algorithm. 
%Next, we will show our branching rules and updates in pricing problem. 

\subsubsection {Branching Rules} 
\label{branchingRule}
Devising good branching rules is an important step in developing a branch-and-price algorithm.
The branching rules may affect the structure of the pricing problem, which may cause the problem 
to become more difficult to solve than the original pricing problem. Therefore, the main 
objective is to find branching rules that eliminate as much of the integer infeasible region 
as possible but that do not increase the difficulty of the pricing problem too much. We apply 
the following branching rules in our algorithm.
\begin {itemize}
\item [{\bf Rule 1}]

{\bf Branching on the facility location variables:} At any node of the branch-and-price tree, 
if the corresponding LP solution has at least one fractional facility location variable, then
we select one of the fractional facility location variables, say $T_j$, as our branching
variable. 
\begin {itemize}
\item On the left branch, we set $T_j=1$, meaning that the corresponding facility is forced 
  to be open. The pricing problem for the corresponding node does not change.
\item On the right branch, we set $T_j=0$, which forces facility j to be closed. 
  If facility j is closed, then no pairings associated with facility j are allowed in the 
  solution. Thus, we fix the variables corresponding to pairings of facility j to zero. The structure 
  of the pricing problem does not change. In fact, we do not solve the pricing problem for facility j, 
  since it is fixed closed.
\end{itemize}

\item [{\bf Rule 2}]

  {\bf Branching on the total number of vehicles:} By definition, each pairing 
  corresponds to one vehicle. Also, vehicles (pairings) belong to a single facility. Thus, in 
  any integer solution, the number of vehicles (pairings) used for any facility must be integer. 
  If we define $nV_j$ to be the number of vehicles used at facility j, then we have the following 
  set of equalities: 
  \begin{eqnarray}
    \label{n-equal}
   nV_j =\sum_{p \in P_j} Z_{jp} \qquad \forall j \in M 
  \end{eqnarray}      
  
%add how we branch on the new variables

  At any node of the branch-and-price tree, if any of the $nV_j$ variables have fractional values, then 
  we select   one as our branching variable. If the selected variable $nV_j$ has the fractional value 
  $a$, then  we create two new branches:
  \begin{eqnarray*}
   nV_j & \leq & \lfloor a \rfloor \qquad  \mbox {(on the left branch)}\\
   nV_j & \geq & \lceil a \rceil    \qquad \mbox {(on the right branch)}
  \end{eqnarray*}    
  
  To do this branching, we introduce the variables $nV_j$, $\forall j\in M$, and add the equations 
  (\ref{n-equal}) to our formulation. Branching on the $nV_j$ variables instead of implicitly  branching
  on $\sum_{p \in P_j}Z_{jp}$ is an example of explicit constraint branching, an idea developed by 
  \citet{Appleget} to improve the performance of the branch-and-bound algorithms for mixed integer
  linear problems. One benefit of these new variables is that it is easier to see the necessary 
  updates to the pricing problem. Let $v_j$ be the dual variable associated with constraint 
  (\ref{n-equal}). Then, the reduced cost of the variable $Z_{jp}$  changes to the following:
  \begin{eqnarray}
    \label{new-red-cost}
    \hat{C}_{jp}=C_{jp}-\sum_{i\in N}a_{ipj}.\pi_{i}+\sum_{i \in N}a_{ipj}.d_{i}.\mu_{j}+\sum_{i \in N}a_{ipj}.\sigma_{ji}-v_j
  \end{eqnarray} 
  In equation (\ref{new-red-cost}), $v_j$ can be considered as a fixed cost for facility j that is 
  shared with every customer node in the pairing p. We do not need to change our network structure or the pricing 
  problem structure to incorporate $v_j$. We calculate the cost of a path in the modified network in same way as 
  before the branching rule and then add the fixed cost $-v_j$ to the cost of the shortest path 
  to find the true reduced cost of the column with the branching rule.    

\item[{\bf Rule 3}]
  {\bf Branching on assignment of a customer to a facility} 
At any node of the branch-and-price tree, if the solution is fractional and Rules 1 and 2 are not 
applicable, then we use rule 3 that is assignment of a customer to a specific facility. 
In a fractional solution, 
if a customer node is served by more than one facility, we select a customer node  $i$ and a facility 
$j$ such that 
 \begin{eqnarray}
\label{branchingrule3-cond}
\sum_{p \in P_j }a_{ipj} Z_{jp}
\end{eqnarray}  
 is fractional. We create two new branches:
  \begin{eqnarray*}
  \sum_{p \in P_j }a_{ipj} Z_{jp} & \geq & 1 \qquad  \mbox {(on the left branch)}\\
   \sum_{p \in P_j }a_{ipj} Z_{jp} & \leq & 0   \qquad \mbox {(on the right branch)}
  \end{eqnarray*}    

In one branch, we force the customer node $i$ to be served by facility $j$. To incorporate this 
assignment, we add the following inequality to the 
RMP as a constraint:
\begin{eqnarray}
  \label{con-rul3}
  \sum_{p \in P_j}a_{ijp}Z_{jp} \geq  1  \ \ (\gamma_{ji})
\end{eqnarray}       
In the pricing problem of facilities other than facility $j$, we simply delete the customer $i$ from the 
network not to generate any pairings including customer $i$. For the pricing problem of facility $j$, 
let $\gamma_{ji}$ be the dual variable assigned to constraint (\ref{con-rul3}). Then, update 
the reduced cost of variables (including customer $i$) for facility $j$ with $\gamma_{ji}$. In
the pricing problem for facility $j$, the cost of an arc into node $i$ is changed to 
  \begin{eqnarray}
\label{redcost-r3}
c_{ki}-\pi_{i}+d_i \cdot \mu_j +\sigma_{ji}-\gamma_{ji}
\end{eqnarray} 
  
In the other branch, we forbid assignment of customer node $i$ to facility $j$. 
To apply the rule, we add the following inequality to the RMP of the current node as a constraint: 
\begin{eqnarray}
  \sum_{p \in P_j }a_{ijp}Z_{jp} \leq  0 \nonumber 
\end{eqnarray}  
In the pricing problem of facility $j$, to create feasible pairings for the branch, 
we remove the node $i$ from the network. The other pricing problems do not change. 

\item[{\bf Rule 4}]
  {\bf Branching on pairs of nodes in the pairing} At any node of the branch-and-price tree,
  if the solution is fractional and Rules 1, 2  and 3 are not applicable, then we use the modified
  \citet*{Ryan} branching rule developed by \citet{Desrochers4}. 

  In the original \citet*{Ryan} branching rule, a pair of set partitioning constraints is 
  selected. In one branch, the rule is to branch on the sum of the variables that satisfy both 
  constraints and, in the other branch, the rule is to branch on the sum of the variables that do 
  not satisfy both constraints simultaneously. In the LRS context, the rule is to select 
  a pair of customer nodes and to branch on the sum of pairings such that the pair of customer  
  nodes appear in the same pairing simultaneously or not. \citet{Desrochers4} modify the 
  Ryan and Foster branching rule by looking only at the variables in which the selected pair of 
  nodes are in consecutive order or not. 
  
  To apply the modified Ryan and Foster branching rule in the LRS context, we select a pair 
  of nodes $(n_1,n_2)$ such that  
  \begin{eqnarray}
    \label{branchingrule4-con}
    0 < \sum_{j \in M}\sum_{\{p \in P_j | (n_1\to n_2) \in p\}} Z_{jp} <1
  \end{eqnarray}       
  where $\{p \in P_j | (n_1\to n_2) \in p\}$ is the set of pairings (of facility $j$)
  in which $n_1$ is followed directly by $n_2$. Based on the structure of the pairings in the LRS 
  problem, if $n_1$ is a customer node, then $n_2$ can be another customer node or it can be the sink; if
  $n_1$ is the source, then $n_2$ can only be a customer node and if $n_1$ is the sink, then $n_2$ can 
  only be a customer
  node.  For example, if we let $i$ and $k$ be two customer nodes, then the $(n_1,n_2)$ pair 
  can be one of the following, $(source,i)$, $(i,k)$, $(i,sink)$ and $(sink,i)$. In one branch, we 
  force the two nodes to be   in consecutive order in any selected pairing and, in the other branch, 
  we forbid any pairings such that the two nodes are in consecutive order. We create two new nodes: 
  \begin{eqnarray}
    & \mbox{Left branch:} & \sum_{j \in M}\sum_{\{p \in P_j | (n_1\to n_2) \in p\}} Z_{jp} \leq  0 \nonumber \\
    & \mbox{Right branch:} & \sum_{j \in M}\sum_{\{p \in P_j | (n_1\to n_2) \in p\}} Z_{jp} \geq  1 \nonumber 
  \end{eqnarray}       

  The left branch specifies that $n_2$ cannot follow $n_1$ in a feasible pairing. The feasible pairings 
  may include $n_1$ or $n_2$ only or may include both of them as long as $n_2$ does not follow $n_1$ directly. 
  In the RMP for the left branch, we fix to zero the previously created pairings in which $n_1$ is 
  followed by $n_2$. In the pricing problem, to create feasible pairings for the left branch, we delete
  the arc connecting node $n_1$ to node $n_2$, so that  no pairing is created in which $n_1$ is followed by $n_2$.
 
  The right branch specifies that either node $n_1$ is followed by node $n_2$ in a pairing or that neither 
  node appears. Pairings that include only one of the nodes, $n_1$ or $n_2$, or that include both nodes
  but $n_2$ does not immediately follow $n_1$ are not feasible in the right branch. Therefore, in the RMP, the 
  variables for previously created infeasible pairings are set to zero. In the pricing problem, to create feasible
  pairings for the right branch, we delete all arcs going into $n_2$ except the arc $(n_1,n_2)$ and
  we delete all arcs going out of $n_1$ except the arc $(n_1,n_2)$. Therefore, if the path visits $n_1$,
  it must visit $n_2$ to be able to reach the sink. To delete an arc in the modified network, we set the cost of 
  the arc to infinity. Other than that, for this branching rule, we do not need to adjust the arc costs. 
\end{itemize}

\subsection {Terminating The Branch-and-Price Tree}
\label{BPAlg-Opt}
After applying the three branching rules to the problem, the  
branch-and-price algorithm terminates with leaf nodes satisfying one of the 
following conditions:
\begin{enumerate}
\item The LP for the node was proved to be infeasible and the node was fathomed.
\item The lower bound of the LP for the node was larger than the best integer feasible 
  solution, so the node was fathomed. 
\item The solution of the LP for the node is integer. Therefore, either the solution is a candidate 
  solution or the node was fathomed because its objective function value was larger than 
  the best integer feasible solution.  
\item The solution of the LP for the node is not integer, but the node is not fathomed 
  because of its objective function value and none of the branching rules can be 
  applied, i.e., all the location variables are integer, the number of vehicles
  for each facility is integer, each customer is served by one facility and there 
are no pairs of nodes that   satisfy condition (\ref{branchingrule4-con}).
\end{enumerate}

In a general branch and bound the algorithm terminates with the leaf nodes all 
satisfying conditions 1, 2 or 3. But in our branch and price tree, 
it is possible to  have leaf nodes satisfying condition 4. 
In this case, we can easily see that the LP solution of the node is equal to
the IP solution of the node. We can find an IP solution by updating the LP solution. 

We can represent a pairing in terms of the nodes or the individual paths included in the 
pairing. Let $(source-1-2-sink-5-7-sink-3-9-4-sink)$ be a feasible pairing visiting 
customer nodes in the set $\{1,2,3,4,5,7,9\}$. We can view this pairing also as a 
combination of the following three paths: $P_1=source-1-2-sink$, $P_2=sink-5-7-sink$ and 
$P_3=sink-3-9-4-sink$. We say that two paths $P_i$ and $P_j$ are {\em node disjoint} if they do not have a 
common customer node, e.g., $P_i\cap P_j =\emptyset$. All of the paths included in a 
feasible pairing must be disjoint from the other because a customer may be visited
at most once. However, in a fractional solution, one or more customer nodes may 
be included in more than one pairing. And, in fact, at least two fractional pairings 
must include either the same path or different but non-disjoint paths. 
Since at the leaf node, we cannot use branching rule 3, we can use 1 facility
in our examples. In an instance with 9 customers and 1 facility, the following solutions might
appear as a fractional solution. Solution 1 includes node disjoint  paths and solution 2 includes
non-disjoint paths.   

\noindent \underline{Solution 1: Same paths:}
\begin{eqnarray*}
0.5 &: & facility_1-1-2-facility_1 \\ 
0.5 &: & facility_1-1-2-facility_1-3-4-5-6-facility_1 \\
0.5 &: & facility_1-7-8-9-facility_1-3-4-5-6-facility_1 \\
0.5 &: & facility_1-7-8-9-facility_1 
\end{eqnarray*}
\underline{Solution 2: Non-disjoint paths:}
\begin{eqnarray*}
0.5 &: & facility_1-1-2-facility_1-3-4-5-6-facility_1 \\
0.5 &: & facility_1-4-6-7-facility_1-8-9-facility_1 \\
0.5 &: & facility_1-5-8-facility_1-7-1-2-3-facility_1-9-facility_1 
\end{eqnarray*}

\begin{lemma} If a leaf node satisfies condition 4, then the fractional 
solution at the node includes 
%pairings that have disjoint paths. Since the solution is fractional, 
at least two pairings that have at least one identical path. And all the paths in 
the pairings are node disjoint from each other (similar to solution 1).
\end{lemma} 
\begin{proof} 
Assume that we have a leaf node satisfying condition 4 and
%and a non-integer solution. 
the fractional solution includes pairings 
that have non-disjoint paths. (The paths, $P_1$ and $P_2$, are non-disjoint
if $(P_1\cap P_2)\neq \emptyset$ and $P_1\neq P_2$). Therefore, there is at least one
arc $(k,i)$ and two paths $P_1$ and $P_2$ such that $i,k \in N$, 
$\{i\} \subseteq (P_1\cap P_2)$ and $(k,i) \in P_1$ but $(k,i) \notin P_2$
and, there is a $k^{'}$ such that $(k^{'},i) \in P_2$.
%Since $P_1$ and $P_2$ are non-disjoint, there should be a $k^{'}$ such that 
%$(k^{'},i) \in P_2$. 
In addition to that, 
$\sum_{j \in M}\sum_{p \in P_j}a_{ijp} Z_{jp}$ must be $1$. Therefore, 
because of the $(k^{'},i)$ link, $0 < \sum_{j \in M}\sum_{\{p \in P_j | (k\to i) \in p\}} Z_{jp} <1$
must be satisfied. This contradicts the assumption that the node satisfies
condition 4. \qed
\end{proof}
%Let $Z_{j1}$ and $Z_{j2}$ be two of the fractional pairings that have at least 
%one non-disjoint pair of paths. Let $P_1 \in Z_{j1}$ and $P_2 \in Z_{j2}$ be 
%the non-disjoint paths, e.g., $(P_1\cap P_2)\neq \emptyset$ and $P_1\neq P_2$.  

% there exists a customer node $i$ that is . Let $P_1 \in Z_{j1}$
%and $P_2 \in Z_{j2}$ such that $\{i\} \subseteq (P_1\cap P_2)$ and
%$P_1\neq P_2$.  

If we have a fractional solution satisfying condition 4, it has a solution similar
to Solution 1 for each open facility. Since each pairing is feasible with respect to
time limit and vehicle capacity, we can delete paths which appear multiple times and
obtain a solution which includes each path (that is in LP solution) exactly once
which is integer.  


%%%%%%%move all these toimplementation details????
%%%%%%%%%%%%%%%%%%%%%%%%%%%%%%%%%%%%%%%%%%%%%%%%%%%%%%%%%%%%%%%%%%%%%%%%%%%%%%%%%%%%%%%%%55
%To implement this branching rule, we created a column pool and path pool. The column pool includes
%every column created with CG algorithm. The path pool includes every path that is included in at 
%least one of the pairings in the column pool. First we branch on the location variables. When 
%there are no fractional location variables, we branch on paths. We construct a network in which 
%each path in the path pool is represented by a node. We add a source node and a sink node to 
%the network. We carry the cost and the total travel time of each path on to the incoming arcs of 
%the node representing that path. Also, each node has a demand that is equal to the total demand 
%of the path corresponding to that node. An adjacency matrix is constructed according to the customers 
%in each path. Suppose that P1=(source-1-2-sink) and P2=(source-2-3-4-sink). There is no arc 
%between P1 and P2 in the constructed network since both paths include the customer node 2. Otherwis,
%if we were to form a pairing with P1 and P2, it would not be integer feasible because of the set 
%partitioning constraint (\ref{spp_con1}). We include arcs from the source to all nodes (paths), 
%arcs from all nodes (paths) to the sink and all arcs between two path nodes that are adjacent according 
%to the adjacency matrix.

%The ESPPRC algorithm is used over this network to find an elementary shortest path from the source to 
%the sink and new columns are created according to the branching rules. Branching rules for the paths 
%are applied with some modifications on the network. For the left branch, we set the travel time 
%between path1 and path2 to infinity. For right branch, we set the travel time on all the arcs going out 
%of path1 and all the arcs going into path2 to infinity except the arc between path1 and path2. 
%With this modification, the structure of the pricing problem does not change but the new pairings 
%satisfy the branching rules. Also, before we process the left or the right branch, we must fix 
%the values of some variables. For example, for the left branch, the pairings in which path2 follows path1
%are set to zero. For the right branch, the pairings in which there is only one of path1 or path2 or 
%both path1 and path2 but in incorrect order (not like $path1\to path2$) are fixed to zero.
     
%In some of our solutions a fractional pairing may include only one path. If this is the case, we may 
%apply our branching strategy to the arc between source and path1 similar to the arc between path1 and path2.

\subsection {2-Step Branch and Price Algorithm}    
\label{sec:2SBPAlgorithm}
In this section, we describe a 2-step branch and price algorithm. Each
step includes solving a branch and price tree. The purpose of the first step is to 
obtain a good upper bound and good columns. The second step is used either to
prove the optimality of the solution or to evaluate the quality of the upper bound. 


\zaupdate{{\bf Create a section for heuristic solutions to LRSP then put the following section} \\
**********************************************8\\
*************************************************
}
  
\subsubsection{Initial Columns and Upper Bound}

To be able to start the first step of the algorithm, we generate some initial columns 
and an upper bound for the LRSP. To do that, we develop a heuristic algorithm 
%called COMBHEUR 
which is a combination of a facility location heuristic, VRP heuristics 
and a bin packing heuristic. We use DROP heuristic to determine a set of facilities 
to be open. Basically, the DROP heuristic starts with all facilities are being open. 
At each iteration, one facility which results the most cost saving is closed until
the problem becomes infeasible. Total cost which includes the facility fixed 
costs, vehicle routing costs and vehicle fixed costs is calculated at every 
iteration of DROP heuristic for a given set of open facilities. 

In the total cost calculation procedure (TCOST), the first step is for any open set of facilities 
to determine an assignment of customers based on the distance and facility capacity. 
For each open facility and given customer assignment, using a VRP 
heuristic vehicle routes are constructed to serve the demands of customers assigned to
the facility. After designing the routes, a bin packing heuristic (best fit decreasing)
is used to combine the routes considering the time limit for the vehicles. This
decides the assignment of vehicles to routes as well as the number of vehicles, 
therefore the fixed cost of vehicles. The total cost calculation is repeated for each 
of the following VRP heuristics: Clark and Wright, nearest neighbor, nearest insertion 
and sweep heuristic. The minimum total cost obtained from the total cost calculation
procedure using one of the VRP heuristics is returned to the DROP heuristic to 
determine the facility to be closed. 

When the DROP heuristic terminates, it
returns a set of open facilities and the total cost of the solution which includes 
the routing and scheduling information as well as the facility fixed costs.
We add the columns which represent the solution from the combined heuristic
to the RMP and update the upper bound. To balance the number of columns added 
for each facility, for the facilities that are fixed to be closed in the heuristic 
solution, we determine a set of customers that can be served from the facility within 
the time limit and use VRP heuristics and the bin packing heuristic in combination to 
determine feasible pairings. 


\subsubsection {Step 1}
Step 1 starts with the upper bound and the RMP that includes the columns generated 
using the heuristic described above. Using the heuristic ESPPRC algorithm (with label 
limit at each node) as a pricing problem and the branching rules described in the 
section \ref{branchingRule}, branch and price algorithm is run for a limited time. When
the algorithm terminates or the time limit is exceeded, the algorithm results
a feasible solution which is an upper bound for the LRSP and a set of feasible pairings.  
 
\subsubsection{Step 2}    
%Step 2 uses the upper bound and the pairings obtained at step 2 as an input. 
Step 2 starts with the upper bound and the RMP that includes the pairings obtained
at step 1. Using the combination of heuristic and exact ESPPRC algorithm as a 
pricing problem the branch and price algorithm is run. Pricing problem, first
uses the heuristic ESPPRC, if the heuristic algorithm cannot find any columns with 
negative reduced cost, the label limit is increased. After some point, if the 
heuristic ESPPRC cannot find a column, exact pricing problem is solved. Together
with the branching rules exact pricing problem guarantees that the solution 
found at the end of the algorithm is optimal solution for the LRSP. Figure \ref{2StepBP}
shows the flow of 2-Step Branch and Price Algorithm. 
\begin{figure}[h]
  \centering
  \includegraphics[scale=0.7]{figs/bpalgorithm.eps}
  %\fbox{ \includegraphics[scale=0.5]{figure/pair2.eps}}
  \caption{2-Step Branch and Price Algorithm} \label{2StepBP}
\end{figure}
